% from ega4-15-fr.tex
\oldpage[IV-3]{224}
Les fibres $f^{-1}(y')$ de $f$ (pour tout $y'\in Y$) ne changeant pas lorsque
l'on remplace $f$ par $f_{\mathrm{red}}$, on peut supposer que $X$ et $Y$ sont
réduits (2.4.3, (vi)), donc $Y$ intègre~; on posera
$K=\mathbf{k}(\eta)$. On peut en outre remplacer $f$ par la restriction
$\operatorname{Supp}(\mathcal F)\to Y$ de $f$ au sous-préschéma fermé réduit de
$X$ ayant $\operatorname{Supp}(\mathcal F)$ comme espace sous-jacent, donc se
borner au cas où $X=\operatorname{Supp}(\mathcal F)$. Enfin, si $y=\eta$, il n'y
a qu'une seule composante irréductible $X_i$ de $X$ contenant $z$
($\mathbf{0}_{\mathrm I}$, 2.1.8) et on a $z_i=z$, ce qui rend (15.1.1.1) et
(15.1.1.2) triviales (sans hypothèses sur $f$). On peut donc se borner au cas
$y\ne\eta$~; il résulte alors de l'hypothèse et de (1.10.3) que les $z_i$
appartiennent à $f^{-1}(\eta)$, les seconds membres de (15.1.1.1) et
(15.1.1.2) étant donc définis. Soit $Z_i$ le sous-préschéma fermé réduit de
$f^{-1}(\eta)$ ayant pour espace sous-jacent $X_i\cap f^{-1}(\eta)$~; on sait
(4.6.6) qu'il existe une extension finie radicielle $K'$ de $K$ telle que, pour
tout $i$, le $K'$-préschéma $(Z_i\otimes_K K')_{\mathrm{red}}$ soit
géométriquement réduit sur $K'$. En outre, le morphisme projection
$f^{-1}(\eta)\otimes_K K'\to f^{-1}(\eta)$ est fini, dominant et radiciel
(\textbf{I}, 3.5.7), donc est un homéomorphisme (2.4.5). Si $z_i'$ est l'unique
point de $f^{-1}(\eta)\otimes_K K'$ au-dessus de $z_i$, il résulte de (4.7.9) et
(4.7.5) que l'on a
\[
\tag{15.1.1.3}
\label{IV.15.1.1.3-fr}
  \lambda_{z_i}(\mathcal F_\eta)
  =\lambda_{z_i'}(\mathcal F_\eta\otimes_K K')
  =\operatorname{long}((\mathcal F_\eta\otimes_K K')_{z_i'}).
\]

Soit $V$ un anneau de valuation discrète, de corps de fractions $K'$, dominant
$\mathcal O_y$ (\textbf{II}, 7.1.7)~; posons $Y'=\operatorname{Spec}(V)$,
$X'=X\times_Y Y'$, $\mathcal F'=\mathcal F\otimes_Y Y'$,
$f'=f_{(Y')}:X'\to Y'$~; si $g:Y'\to Y$ est le morphisme structural, $\eta'$ le
point générique de $Y'$ et $y'$ son point fermé, on a $g(y')=y$, $g(\eta')=\eta$~;
le morphisme $\operatorname{Spec}(\mathbf{k}(y'))\to
\operatorname{Spec}(\mathbf{k}(y))$ étant fidèlement plat, il en est de même de
la projection $g':f'^{-1}(y')\to f^{-1}(y)$ (2.2.13), donc (2.3.4), il y a un
point générique $z'$ d'une composante irréductible de $f'^{-1}(y')$ tel que
$g'(z')=z$. Par construction, on a $\mathbf{k}(\eta')=K'$, donc
$f'^{-1}(\eta')=f^{-1}(\eta)\otimes_K K'$~; d'autre part, si l'on pose
$X_i'=X_i\times_Y Y'$, la restriction $X_i'\to Y'$ de $f'$ est un morphisme
universellement ouvert au point $z'$, donc (1.10.3) il existe un point
$z_i''\in f'^{-1}(\eta')$ qui est une générisation de $z'$, et l'on peut
évidemment supposer que $z_i''$ est point générique d'une composante
irréductible de $X_i'\cap f'^{-1}(\eta')$~; comme la projection de
$X_i'\cap f'^{-1}(\eta')$ dans $f^{-1}(\eta)$ est $Z_i$, on a $z_i''=z_i'$.
En vertu de (4.7.9) et (4.7.5), on a
\[
\tag{15.1.1.4}
\label{IV.15.1.1.4-fr}
  \lambda_z(\mathcal F_y)=\lambda_{z'}(\mathcal F'_{y'})
  \geqslant\operatorname{long}((\mathcal F'_{y'})_{z'})
\]
et il résulte donc de cette inégalité et de (15.1.1.3) qu'il suffit, pour établir
(15.1.1.1), de démontrer l'inégalité
\[
\tag{15.1.1.5}
\label{IV.15.1.1.5-fr}
  \operatorname{long}((\mathcal F'_{y'})_{z'})\geqslant
  \sum_{i=1}^{n}\operatorname{long}((\mathcal F'_{\eta'})_{z_i'}).
\]

Considérons maintenant la seconde inégalité (15.1.1.2)~; nous utiliserons le
lemme suivant~:

\begin{lemma}[15.1.1.6]
\label{IV.15.1.1.6-fr}
Soient $A$ un anneau local régulier qui n'est pas un corps, $K$ son corps des
fractions, $k$ son corps résiduel. Il existe un anneau de valuation discrète $V$
dominant $A$, ayant $K$ pour corps des fractions et dont le corps résiduel est
une extension transcendante pure de $k$.
\end{lemma}

